\section{Introducción y Público Objetivo}

\subsection{El Problema a Resolver}
En la actualidad, la adquisición de boletos para eventos en vivo suele ser un proceso engorroso o propenso a fraudes en la reventa. Los usuarios se enfrentan frecuentemente a los siguientes inconvenientes:
\begin{itemize}
    \item \textbf{Transferir entradas:} Imposibilidad de transferir boletos de forma segura a amigos o familiares sin incurrir en cobros abusivos de plataformas intermediarias.
    \item \textbf{Cancelación de reservas:} Falta de mecanismos rápidos para cancelar boletos cuando surge un imprevisto, impidiendo liberar el asiento a otros interesados en tiempo real.
    \item \textbf{Acceso offline y seguridad:} La necesidad de contar con códigos de barras estáticos o PDFs que son fáciles de duplicar y no brindan una garantía de seguridad integrada al dispositivo.
\end{itemize}

\subsection{Público Objetivo}
El público objetivo de \textbf{EventTix} se clasifica en dos segmentos fundamentales:
\begin{itemize}
    \item \textbf{Usuarios Finales (Espectadores):} Personas de entre 16 y 60 años habituadas al consumo digital, que asisten a conciertos, obras teatrales y eventos deportivos, y exigen transacciones fluidas y seguras.
    \item \textbf{Organizadores de Eventos:} Promotores de entretenimiento que requieren visibilidad en tiempo real sobre el estado de la venta de entradas, cancelaciones y ocupación de asientos en el recinto.
\end{itemize}

\subsection{Credenciales de Usuarios de Prueba (Seeds)}
Para facilitar la evaluación de la funcionalidad en tiempo real y flujos P2P, la base de datos se encuentra sembrada con las siguientes cuentas de prueba:
\begin{itemize}
    \item \textbf{Usuario Administrador:}
    \begin{itemize}
        \item Correo electrónico: \texttt{admin@eventtix.com}
        \item Contraseña: \texttt{admin123}
    \end{itemize}
    \item \textbf{Usuario Secundario:}
    \begin{itemize}
        \item Correo electrónico: \texttt{user@eventtix.com}
        \item Contraseña: \texttt{user123}
    \end{itemize}
\end{itemize}

